\subsubsection{Limitations}

Several limitations constrain the generalizability of our findings.

\textbf{Reward design scope.}
All rewards are designed specifically for GSM8K's answer format (\texttt{\#\#\#\# <number>}).
Format compliance and numeric proximity both depend on this structure and would not transfer to tasks with non-numeric or free-form answers without redesign.

\textbf{Reward hacking risk.}
The format compliance reward incentivizes producing the \texttt{\#\#\#\#} marker regardless of reasoning quality.
Without KL regularization, the model could learn to produce well-formatted but poorly-reasoned responses.
Similarly, numeric proximity could incentivize guessing numbers close to common answer magnitudes rather than performing actual computation.

\textbf{Scale.}
We evaluate on a single SFT checkpoint, a single dataset (GSM8K), and a single RL algorithm (REINFORCE).
Results may not generalize to stronger base models, harder math benchmarks, or more sophisticated RL algorithms such as PPO or GRPO with KL penalties.

\textbf{Reward composition.}
Rewards are summed with equal weight.
We do not explore reward weighting, scheduling, or curriculum strategies that might improve the Comb.\ All run's performance.

\textbf{Statistical significance.}
We report single-run results without confidence intervals or multiple seeds.
Differences below 2 percentage points should be interpreted with caution, and even larger differences may not be robust across random seeds.

\textbf{Mistake taxonomy.}
Our rule-based taxonomy (Section 4.2) uses simple heuristics (character count thresholds, regex patterns, relative error cutoffs) that may misclassify edge cases.
For example, a response with correct reasoning but a transcription error in the final line would be classified as ``arithmetic error'' or ``large error'' rather than a more nuanced category.
